
\begin{frame}{Что такое POK?}
    \begin{itemize}
        \item \textbf{Ядро разделения} (Separation Kernel) для критических по безопасности систем.
        \item Полное соответствие стандарту \textbf{ARINC 653} и архитектуре \textbf{MILS}.
        \item Основная цель: гарантировать, что сбой в одной подсистеме не распространится на другие.
        \item Два типа изоляции:
        \begin{itemize}
            \item \textbf{Пространственная} -- разделы не могут обращаться к чужой памяти.
            \item \textbf{Временная} -- каждому разделу выделен фиксированный квант времени, который невозможно нарушить.
        \end{itemize}
        \item Минималистичное ядро: >90\% кода легко покрывается тестами и нацелено на формальную верификацию.
    \end{itemize}
\end{frame}

\begin{frame}{Двухуровневая архитектура POK}
    \begin{itemize}
        \item \textbf{Ядро}:
        \begin{itemize}
            \item Управление аппаратурой.
            \item Таймеры, переключение разделов, доставка прерываний.
            \item Минимум системных вызовов (\textit{hypercall}).
        \end{itemize}
        \item \textbf{Среда выполнения Runtime/libpok}:
        \begin{itemize}
            \item Библиотека, компонуемая с каждым разделом.
            \item API: потоки ARINC-653, аллокаторы, коммуникации.
            \item Драйверы устройств работают именно в разделах, а не в ядре.
        \end{itemize}
    \end{itemize}
\end{frame}

\begin{frame}{Процесс сборки и загрузки системы}
    \begin{itemize}
        \item Каждый раздел и само ядро компилируются независимо в отдельные ELF-разделы.
        \item Все ELF-разделы упаковываются в единый архив.
        \item При загрузке:
        \begin{enumerate}
            \item Загрузчик запускает ядро POK.
            \item Ядро извлекает каждый ELF.
            \item Каждый раздел загружается в свой изолированный сегмент памяти, защищённый MPU (Memory Protection Unit)/MMU (Memory Management Unit).
        \end{enumerate}
        \item Любое нарушение вызывает отказ.
    \end{itemize}
\end{frame}

\begin{frame}{Временная изоляция}
    \begin{itemize}
        \item Каждому разделу назначается фиксированный временной слот и период активации.
        \item Планировщик ядра работает по циклической таблице.
        \item Ни один раздел не может занять больше времени, гарантируется аппаратным таймером.
        \item Это исключает влияние одного раздела на временные характеристики других.
    \end{itemize}
\end{frame}

\begin{frame}{Коммуникации между разделами}
    \begin{itemize}
        \item Взаимодействие только через \textbf{порты ARINC 653}, настраиваемые на этапе конфигурации.
        \item Два типа портов:
        \begin{itemize}
            \item \textbf{Sampling Port} — хранит последнее значение, перезаписывается.
            \item \textbf{Queuing Port} — очередь сообщений (FIFO).
        \end{itemize}
        \item Каналы связи строго статически определены, не могут быть изменены во время выполнения.
        \item Ядро проверяет права доступа к портам — раздел может писать/читать только разрешённые порты.
    \end{itemize}
\end{frame}

\begin{frame}{AADL и Ocarina}
    \begin{itemize}
        \item Архитектура системы описывается на языке \textbf{AADL} (Architecture Analysis \& Design Language):
        \begin{itemize}
            \item Процессорные платы, память, разделы, потоки, каналы связи.
        \end{itemize}
        \item Инструмент \textbf{Ocarina} автоматически генерирует:
        \begin{itemize}
            \item Конфигурационные файлы POK (XML-описание разделов и портов).
            \item Код инициализации и коммуникацию между разделами.
        \end{itemize}
        \item Сокращает ручное кодирование, минимизирует ошибки и упрощает сертификацию.
    \end{itemize}
\end{frame}

\begin{frame}{Безопасность: уязвимость CVE-2018 (Airbus CERT)}
    \begin{alertblock}{Найденная проблема}
        Airbus Group CERT обнаружил, что недостаточная проверка указателей в системных вызовах POK позволяла вредоносному коду внутри раздела читать/писать память за его пределами.
    \end{alertblock}
    \begin{itemize}
        \item Уязвимость подчёркивает сложность реализации абсолютной изоляции даже в минималистичном ядре.
        \item После исправления: усилена валидация всех указателей, передаваемых из разделов.
        \item Демонстрирует необходимость \textbf{формальной верификации} как кода ядра, так и самих проверок.
    \end{itemize}
    \begin{center}
        \fbox{\textbf{Урок:} Даже в разделенном ядре ошибка валидации может разрушить всю изоляцию.}
    \end{center}
\end{frame}

\begin{frame}{Применение в реальных системах}
    \begin{itemize}
        \item \textbf{Промышленность:} Форк POK -- \textbf{JetOS} (ИСП РАН) -- используется в коммерческих авиалайнерах.
        \item \textbf{Научные исследования:} Доверенное ПО для квадрокоптеров на базе POK, верифицированное с AADL.
        \item \textbf{Обучение:} Платформа для изучения архитектур реального времени и безопасных систем.
    \end{itemize}
\end{frame}

\begin{frame}{Резюме}
    \begin{columns}[T]
        \column{0.5\textwidth}
        \textbf{Сильные стороны}
        \begin{itemize}
            \item Жёсткая пространственная и временная изоляция.
            \item Соответствие ARINC 653 / MILS.
            \item Минимальный код ядра.
            \item AADL/Ocarina.
            \item Открытая архитектура, пригодная для верификации.
        \end{itemize}
        \column{0.5\textwidth}
        \textbf{Вызовы и риски}
        \begin{itemize}
            \item Необходимость освоения AADL и инструментария.
            \item Неисправленные уязвимости (пример CVE).
            \item Требуется дальнейшая формальная верификация для наивысших уровней гарантий.
        \end{itemize}
    \end{columns}
    \vfill
    \begin{center}
        \textbf{POK -- пример правильной архитектуры для систем, где критически важна безопасность.}
    \end{center}
\end{frame}